\song{Mladičká básnířka (Jaromír Nohavica)}

\ve{}
\ch{C}Mladičká básnířka s \ch{Emi}korálky nad kotníky\ch{Ami G}\\
\ch{C}bouchala na dvířka \ch{Emi}paláce poetiky\ch{Ami G}\\
s někým se \ch{C}vyspala někomu \ch{G}nedala\\
láska jako \ch{Ami}hobby \ch{Cmi}pak o tom napsala \ch{G}blues\\
na čtyři \ch{C}doby, hoho\ch{Ami}ho hoho\ch{Emi}ho ho \ch{Ami}ho ho ho ho\ch{G}

\ve{}
Své srdce skloňovala podle vzoru Ferlinghetti\\
ve vzduchu nechávala viset vždy jen půlku věty\\
plná tragiky plná mystiky plná splínu\\
pak jí to otiskli v jednom magazínu

\ve{}
Bývala viděna v malém baru u Rozhlasu\\
od sebe kolena a cizí ruka kolem pasu\\
trochu se napila trochu se opila\\
na účet redaktora za týden nato byla hvězdou Mikrofóra

\ve{}
Pod paží nosila rozepsané rukopisy\\
ráno se budila vedle záchodové mísy\\
životem potřísněná múzou políbená\\
plná zázraků a pak ji vyhodili z gymplu\\
a hned nato i z baráku 

\ve{}
Ve třetím měsíci dostala chuť na jahody\\
ale básníci-tatíci nepomýšlej na rozvody\\
cítila u srdce jak po ní přešla železná bota\\
pak o tom napsala sonnet a ten byl ze života

\re{(6.)}
A jednou v pondělí přišla do klubu na koleje\\
a hlasem nesmělým prosila o text Darmoděje\\
a jak tak psala, náhle se dala tichounce do pláče\\
a inkoustové slzy kapaly na její mrkváče

\re{(7.)}
Jó mladé básnířky\\
vy mladé básnířky\\
ach mladé básnířky

\esong{}